%% 
%% Copyright 2019-2024 Elsevier Ltd
%% 
%% Complex Article Service (CAS) LaTeX template
%% Formatted for Online Social Networks and Media (OSNEM)
%% 

\documentclass[a4paper,fleqn]{cas-dc}

\usepackage[authoryear]{natbib}
\usepackage{amsmath,amssymb,amsfonts}
\usepackage{graphicx}
\usepackage{booktabs}
\usepackage{microtype}
\usepackage{url}
%% \usepackage{doi} % removed to eliminate doi:doi: duplication in cas-model2-names

%%% Author definitions
\def\tsc#1{\csdef{#1}{\textsc{\lowercase{#1}}\xspace}}
\tsc{WGM}
\tsc{QE}

\begin{document}
\let\WriteBookmarks\relax
\renewcommand{\topfraction}{0.95}
\renewcommand{\bottomfraction}{0.95}
\renewcommand{\textfraction}{0.05}
\renewcommand{\floatpagefraction}{0.40}
\renewcommand{\dbltopfraction}{0.95}
\renewcommand{\dblfloatpagefraction}{0.40}
\setcounter{topnumber}{4}
\setcounter{bottomnumber}{4}
\setcounter{totalnumber}{8}
\setcounter{dbltopnumber}{4}

\shorttitle{Opinion--conviction decoupling explains late-campaign electoral swings}
\shortauthors{C. Rennó-Costa}

\title [mode = title]{Opinion--conviction decoupling explains late-campaign electoral swings and how socio-technical counter-strategies neutralize them}

\author[1]{César Rennó-Costa}[orcid=0000-0002-3932-2327]
\cormark[1]
\ead{cesar@imd.ufrn.br}

\affiliation[1]{organization={Digital Metropolis Institute and Brain Institute, Federal University of Rio Grande do Norte (UFRN)},
                addressline={Av. Sen. Salgado Filho, 3000, Lagoa Nova},
                city={Natal},
                postcode={59078-970},
                state={RN},
                country={Brazil}}

\cortext[1]{Corresponding author}

\begin{highlights}
\item Two-timescale model captures fast opinion shifts away from deep convictions.
\item Negativity hysteresis (3.5x kinetic ratio) favors late disinformation blitzes.
\item 2026 Brazilian runoff reversal serves as an empirical quasi-natural experiment.
\item Janonismo counter-tactics systematically neutralize algorithmic outrage feeds.
\end{highlights}

\begin{abstract}
Voters' core political convictions exhibit strong temporal stability, yet democratic elections can experience abrupt late-campaign volatility that confounds pre-election polling. While single-round polling discrepancies have occurred across major democracies, experiencing consecutive late-campaign reversals in opposite directions across sequential rounds is an empirical anomaly. This study analyzes this two-round divergence using an empirical digital twin of the 2026 Brazilian presidential election, where an eve-of-election conservative surge was followed by a runoff reversal. Coupling a census-calibrated urban multiplex network, broadcast media, and algorithmic platforms, the model reproduces longitudinal survey trajectories and official ballot totals ($R^2 = 0.9955$, root-mean-square error $0.85$ percentage points). I show that fast-opinion elasticity temporarily decouples expressed voting choices from stable convictions. Amplified by asymmetric candidate aversion ($3.5\times$ faster accumulation than decay), this decoupling enables dark-social cascades to overturn leads during polling blackouts, while coordinated socio-technical counter-strategies can rapidly neutralize these swings.
\end{abstract}

\begin{keywords}
Opinion dynamics \sep Multiplex networks \sep Algorithmic curation \sep Disinformation cascades \sep Electoral volatility \sep Janonismo Cultural
\end{keywords}

\maketitle

\section{Introduction}
\label{sec:intro}

A central tenet of political behavior is that citizens' core ideological convictions exhibit high temporal inertia, sustained by durable socioeconomic identities, generational socialization, and stable partisan attachments \citep{converse1964nature, alwin1991aging, zaller1992nature}. Over multi-year horizons, deep-seated political values change slowly, acting as psychological anchors that buffer democratic electorates against transient informational shocks \citep{kiley2020measuring, howe2017attitude}. Yet major democratic contests have repeatedly exhibited abrupt, late-breaking swings that defy pre-election survey projections despite the documented stability of underlying ideological convictions \citep{tucker2018social, lorenz2019accelerating, jennings2018election}. Notable precedents include late-stage shifts observed in the 2016 United Kingdom European Union membership referendum and the 2016 United States presidential election \citep{kennedy2018evaluation, jennings2018election}, as well as Brazil's 2018 presidential surge \citep{nicolau2020brasil} and the 2022 Brazilian first round, where incumbent Jair Bolsonaro outperformed eve-of-election polling estimates by over five percentage points \citep{tse2022resultados}. However, while these historical contests involved single late-campaign polling misses in one-round or decisive general elections, experiencing two consecutive late-campaign reversals in opposite directions across sequential rounds of the same presidential contest represents an unusual empirical anomaly.

Deciphering the mechanisms of late-campaign volatility requires moving beyond single-channel models of public opinion. Real-world political deliberation unfolds across three distinct, dynamically coupled communication substrates \citep{boccaletti2014structure, kivela2014multilayer}. First, face-to-face interpersonal interactions across everyday urban environments (households, workplaces, schools, religious congregations, friendship circles, and transit spaces) provide localized social reinforcement and cross-cutting deliberative friction \citep{granovetter1973strength, mcpherson2001birds}. Adapted from census-calibrated multi-layer urban agent-based models developed in computational epidemiology \citep{lopes2022measuring}, realistic sociopolitical digital twins must capture these structured contact environments to represent interpersonal contagion across city scales. Second, traditional mass media institutions (broadcast television, radio, and investigative print journalism) synchronize collective attention and establish professional editorial gatekeeping \citep{mccombs1972agenda, gentzkow2006media}. Third, modern algorithmic social platforms and dark-social encrypted messaging networks curate personal discovery feeds via engagement-maximizing recommendation engines while enabling frictionless peer-to-peer forwarding \citep{cinelli2021echo, vosoughi2018spread, cesarino2020posverdade}. 

A critical theoretical gap concerns how these three communicative substrates interact under acute, time-compressed informational offensives. Traditional campaign institutions rely on statutory broadcast airtime, legacy press conferences, and judicial fact-checking petitions \citep{gentzkow2006media, chadwick2017hybrid}. However, empirical evidence indicates that these defensive mechanisms suffer from a wide separation of operational timescales \citep{downs1972up, harder2017intermedia}; while institutional newsrooms and regulatory courts deliberate over days or weeks, algorithmic social networks propagate moral-emotional cascades in minutes \citep{brady2017emotion, vosoughi2018spread}. Furthermore, contemporary feed recommendation engines optimize engagement symmetrically for both positive affirmation and negative outrage \citep{crockett2017moral, bail2018exposure}. Consequently, defensive debunking and quote-dunking paradoxically supply the engagement signals curation algorithms need to amplify opponent disinformation into unaligned discovery feeds.

The 2026 Brazilian presidential election provides a quasi-natural experiment to investigate these dynamics. Unlike previous contests characterized by a single polling miss, Brazil in 2026 exhibited a genuine double polling anomaly, where tracking surveys failed to capture an acute election-eve conservative surge in the first round and subsequently missed a progressive counter-surge in the second round. Pitting incumbent president Luiz Inácio Lula da Silva against Senator Flávio Bolsonaro, the contest took place in a mobile-first democracy where statutory television advertising and digital networks operated across two distinct regulatory regimes. Throughout the first round, the incumbent coalition maintained a broadcast-centric strategy backed by a $+27.3\%$ advantage in statutory television airtime. Despite steady polling leads over fifty days, the incumbent was overtaken on election eve by an acute digital shock that lifted Flávio Bolsonaro into first place with $47.03\%$ of valid votes, confounding eve-of-election survey projections. In the ensuing runoff, statutory broadcast airtime was equalized to strict 50/50 parity by law, removing broadcast advantages and isolating digital campaign strategy as the primary differentiator. Faced with an initial deficit that tracking polls projected would solidify into defeat, the progressive campaign executed a strategic pivot by adopting the counter-tactical doctrine of \emph{Janonismo Cultural}, articulated by Federal Deputy André Janones \citep{janones2023janonismo}, deploying socio-technical interventions designed to starve rival outrage, concentrate viral bandwidth, and bypass open feeds through dark-social messaging.

To investigate these mechanisms, I construct a computational digital twin of the 2026 Brazilian electorate to evaluate the microscopic drivers of late-campaign electoral shocks and the efficacy of socio-technical counter-tactics. Calibrated to national census demography, detailed tracking surveys, and empirical platform telemetry, the agent-based model couples a five-layer multiplex interpersonal network, an institutional mass media ecosystem, and four algorithmic social platforms across a four-dimensional concept hypercube. By formalizing a two-timescale separation between fast-moving expressed opinions $\mathbf{x}_i(t)$ and slow-crystallizing convictions $\mathbf{s}_i(t)$, this framework captures how transient digital perturbations decouple public opinion from foundational beliefs. The same fast-opinion elasticity $\mathbf{x}_i(t)$ that leaves electorates vulnerable to acute late-campaign shocks also explains why coordinated socio-technical counter-strategies can rapidly neutralize those swings within a short runoff window without needing to reshape slow-moving deep convictions $\mathbf{s}_i(t)$. Precisely because 2026 Brazil exhibited both an initial late-stage conservative surge in Round~1 and a subsequent progressive reversal in Round~2, the contest provides an empirical setting to evaluate how the same opinion-elasticity mechanisms can both destabilize and restore an electoral majority. Through large-scale Monte Carlo ensembles ($N = 100{,}000$ citizens across 100 independent realizations), I isolate the psychological and topological drivers that enable late-stage disinformation blitzes to overturn established majorities during polling blackouts, and evaluate how structured algorithmic counter-offensives can mitigate late-stage erosion to restore electoral resilience.

\section{Related Work}
\label{sec:related_work}

This investigation bridges three previously disconnected literatures in computational social science, network physics, and behavioral political science.

\subsection{Multiplex Social Contact Networks}
Classical models of opinion dynamics typically operate on single, unweighted, or static graphs \citep{castellano2009statistical, hegselmann2002opinion}. However, empirical research in complex networks emphasizes that human interaction unfolds across multiplex topologies \citep{boccaletti2014structure, kivela2014multilayer}. In urban systems, social contagion occurs simultaneously across physical contact structures (households, schools, workplaces, and public transit) and media layers \citep{lopes2022measuring}. Distinct contact environments exhibit starkly different structural properties; households and workplaces show high clustering and strong ties \citep{granovetter1973strength}, while transit environments provide random weak ties that accelerate long-range diffusion. Capturing these heterogeneous contact topologies is essential to understanding why local peer clusters often resist external communicative shocks \citep{mcpherson2001birds}.

\subsection{Two-Timescale Bounded-Confidence Dynamics}
Mathematical models of opinion formation have evolved from binary voting models to continuous bounded-confidence frameworks, notably the Deffuant--Weisbuch \citep{deffuant2000mixing} and Hegselmann--Krause \citep{hegselmann2002opinion} models, as well as the Friedkin--Johnsen social power formulation \citep{friedkin1990social, friedkin2016network}. While these models capture peer consensus and polarization, they often treat expressed opinions as monolithic. In cognitive psychology, however, a fundamental distinction exists between top-of-the-head public opinions, which are volatile and sensitive to immediate framing \citep{zaller1992nature}, and deep ideological convictions, which are crystallized through lifelong socialization \citep{converse1964nature, howe2017attitude, kiley2020measuring}. Although social judgment theory posited latitudes of acceptance and rejection around cognitive anchors \citep{sherif1961social}, dynamical models have rarely formalized this coupling through an explicit two-timescale elastic tether.

\subsection{Algorithmic Curation, Negativity Bias, and Encrypted Cascades}
Modern electoral campaigns are increasingly mediated by proprietary algorithmic recommendation systems \citep{cinelli2021echo, tucker2018social}. Commercial discovery algorithms optimize engagement metrics by treating positive affirmation and hostile outrage symmetrically, amplifying moral-emotional language along the diagonal of controversy \citep{brady2017emotion, crockett2017moral, bail2018exposure}. Concurrently, cognitive psychology establishes a pervasive negativity bias whereby humans process threatening or aversive information more rapidly and retain it longer than positive information \citep{rozin2001negativity, baumeister2001bad}. In electoral contexts, candidate aversion triggers stronger voting mobilization and abstention penalties than positive admiration \citep{almeida2019evangelicals, nicolau2020brasil}. Furthermore, the migration of viral content from open discovery feeds into closed, encrypted messaging applications (WhatsApp, Telegram) enables frictionless peer-to-peer cascades that evade editorial fact-checking and statutory polling scrutiny \citep{resende2019whatsapp, garimella2020images, cesarino2020posverdade}.

\section{The Multi-Layer Opinion Dynamics Model}
\label{sec:model}

To formalize the interplay between interpersonal deliberation, broadcast media, and algorithmic platforms, I construct a multi-layer agent-based digital twin.

\subsection{Multidimensional Concept Space}
\label{subsec:concept_space}
Citizens $i \in \{1, \dots, N\}$, political candidates $k \in \{L, F, O\}$, and media outlets $g \in \{1, \dots, M\}$ are embedded in a $D$-dimensional continuous concept hypercube $\Omega = [-1, 1]^D$ ($D = 4$). In the Brazilian context, coordinate axes correspond to the primary structural cleavages, consisting of fiscal policy ($d=0$), moral, cultural, and religious customs ($d=1$), environmental governance and Amazonian preservation ($d=2$), and public security and firearms regulation ($d=3$).

The normalized metric distance between any two vectors $\mathbf{u}, \mathbf{v} \in \Omega$ is normalized by the maximal Euclidean diameter $2\sqrt{D}$, given by
\begin{equation}
    d(\mathbf{u}, \mathbf{v}) = \frac{1}{\sqrt{D}} \|\mathbf{u} - \mathbf{v}\|_2 \in [0, 2],
    \label{eq:methods_distance}
\end{equation}
such that individual conviction extremity is measured by the metric distance from the ideological origin,
\begin{equation}
    E(\mathbf{u}) \equiv d(\mathbf{u}, \mathbf{0}) = \frac{1}{\sqrt{D}} \|\mathbf{u}\|_2 \in [0, 1].
    \label{eq:methods_extremity}
\end{equation}

Ideological cross-issue constraint is enforced via a row-normalized cognitive coupling matrix $\mathbf{C} = [C_{dd'}] \in \mathbb{R}^{D \times D}$,
\begin{equation}
    \mathbf{C} = \begin{pmatrix}
        0.58 & 0.14 & 0.16 & 0.12 \\
        0.12 & 0.52 & 0.18 & 0.18 \\
        0.14 & 0.18 & 0.50 & 0.18 \\
        0.12 & 0.20 & 0.18 & 0.50
    \end{pmatrix},
    \label{eq:coupling_matrix}
\end{equation}
This matrix embodies the structural constraint observed in empirical legislative roll-call matrices and party manifestos.

\subsection{Urban Multiplex Contact Topology}
\label{subsec:multiplex_network}
Interpersonal interactions unfold across a five-layer multiplex network $\mathcal{M} = \{\mathcal{L}_1, \dots, \mathcal{L}_5\}$ stored in Compressed Sparse Row (CSR) format. Layer 1 ($\mathcal{L}_1$, households) consists of disjoint cliques calibrated to census statistics, with an average of $2.79$ residents per domicile and tie weight $w_1 = 0.36$. Layer 2 ($\mathcal{L}_2$, workplaces and schools) is structured as regular lattices with local clustering ($\langle k^{[2]} \rangle = 10$, tie weight $w_2 = 0.26$). Layer 3 ($\mathcal{L}_3$, religious congregations) is represented by bipartite affiliation graphs linking citizens to local congregations ($\langle k^{[3]} \rangle = 8$, weight $w_3 = 0.14$). Layer 4 ($\mathcal{L}_4$, friendship circles) follows small-world Watts--Strogatz topologies ($p_{\mathrm{ws}} = 0.15$, $\langle k^{[4]} \rangle = 6$, weight $w_4 = 0.14$). Finally, Layer 5 ($\mathcal{L}_5$, transit weak ties) comprises Erd\H{o}s--R\'enyi random graphs representing incidental physical encounters ($\langle k^{[5]} \rangle = 4$, weight $w_5 = 0.10$).

Composite inter-agent weights\linebreak $W_{ij} = \sum_{\ell=1}^5 w_\ell A_{ij}^{[\ell]} \big/ \sum_k \sum_\ell w_\ell A_{ik}^{[\ell]}$ yield a mean interpersonal degree of $\langle k_{\mathrm{comp}} \rangle \approx 30.8$ contacts per citizen. Inter-cohort contact is governed by the empirical age contact matrix $M_{ab}$ across six adult age brackets (16--24, 25--34, 35--44, 45--54, 55--64, 65+).

\subsection{Two-Timescale Opinion--Conviction Dynamics}
\label{subsec:opinion_dynamics}
At each integration step $\Delta t = 0.15$, citizen $i$ engages in peer deliberation with neighbors $j \in \mathcal{N}_i$ satisfying the bounded-confidence condition $d_{ij}(t) \le \epsilon_i + 0.5\,u_i(t)$. Effective peer weights combine Relative Agreement (RA) uncertainty overlap $\phi_{ij}^{\mathrm{RA}}(t)$ and Biased Assimilation (BA) cosine similarity $\phi_{ij}^{\mathrm{BA}}(t)$ as
\begin{equation}
\begin{aligned}
    \tilde{w}_{ij}(t) = \; & W_{ij} \big(0.40 + 0.60\,\phi_{ij}^{\mathrm{RA}}(t)\big) \phi_{ij}^{\mathrm{BA}}(t) \\
    & \times \mathbb{I}\!\big(d_{ij}(t) \le \epsilon_i + 0.5\,u_i(t)\big).
\end{aligned}
\label{eq:methods_peer_weight}
\end{equation}
The interpersonal social shift $\Delta \mathbf{x}_i^{\mathrm{soc}}(t) = \sum_{j} \tilde{w}_{ij}(\mathbf{x}_j - \mathbf{x}_i) / \sum_j \tilde{w}_{ij}$ updates current opinions through the Friedkin--Johnsen formulation,
\begin{equation}
\begin{aligned}
    \mathbf{x}_i^{\mathrm{int}}(t) = \mathrm{clip}_{[-1, 1]^D}\!\Big(& \mathbf{x}_i(t) + \Delta t \big[(1 - \alpha_i)(\mathbf{s}_i(t) - \mathbf{x}_i(t)) \\
    & + \alpha_i \mathbf{C}\Delta \mathbf{x}_i^{\mathrm{soc}}(t)\big] + \boldsymbol{\xi}_i(t)\Big),
\end{aligned}
\label{eq:methods_fj}
\end{equation}
where $\boldsymbol{\xi}_i(t) \sim \mathcal{N}(\mathbf{0}, \sigma_{\xi}^2 \mathbf{I}_D)$ ($\sigma_{\xi} = 0.012$) is intrinsic cognitive noise. When an exogenous factual impulse $\mathbf{b}_{\mathrm{fact}}(t)$ is present, opinions pull toward it with sensitivity $\gamma_i(t) = \gamma_0 (1.30 - 0.72 E_i)(1.18 - 0.48\tilde{a}_i)(0.85 + 0.30\,\mathrm{SES}_i)$, which produces the intermediate opinion $\tilde{\mathbf{x}}_i(t)$.

Opinions cannot depart indefinitely from internal convictions. The maximal allowable departure is defined by an individual tether radius $\Delta_{\max, i}(t) = \Delta_0 (1.25 - 0.55\tilde{a}_i)(1.30 - 0.75 E_i)(0.90 + 0.25\,\mathrm{SES}_i)$. Defining the raw tether excursion as $g_i(t) = \frac{1}{\sqrt{D}}\|\tilde{\mathbf{x}}_i(t) - \mathbf{s}_i(t)\|_2$, radial elastic restraint is enforced via hyperbolic-tangent bounding,
\begin{align}
    g_i^{\mathrm{eff}}(t) &= \Delta_{\max, i}(t) \tanh\!\left(\frac{\|\tilde{\mathbf{x}}_i(t) - \mathbf{s}_i(t)\|_2}{\sqrt{D}\,\Delta_{\max, i}(t)}\right), \label{eq:methods_tether_radius} \\
    \mathbf{x}_i(t+1) &= \mathrm{clip}_{[-1, 1]^D}\!\left(\mathbf{s}_i(t) + \frac{g_i^{\mathrm{eff}}(t)(\tilde{\mathbf{x}}_i - \mathbf{s}_i)}{g_i(t) + 10^{-12}}\right). \label{eq:methods_tether_update}
\end{align}
Convictions adapt over substantially slower timescales through baseline crystallization drift $\eta_{\mathrm{cryst}, i}(t)$ and sustained-tension restructuring $\eta_{\mathrm{sust}, i}(t)$,
\begin{align}
    \eta_{\mathrm{cryst}, i}(t) &= \frac{\eta_0\,(1.25 - 0.75\tilde{a}_i)}{1 + (t/\tau_c)^{1.15}} \nonumber \\
    &\quad \times (1.22 - 0.68 E_i)(1.1 - 0.25\,\mathrm{SES}_i), \label{eq:methods_cryst} \\
    \eta_{\mathrm{sust}, i}(t) &= \frac{0.28\eta_0 (g_i^{\mathrm{eff}})^2}{\Delta_{\max, i}^2} \nonumber \\
    &\quad \times (1.18 - 0.58\tilde{a}_i)(1.20 - 0.62 E_i), \label{eq:methods_sust}
\end{align}
with baseline drift $\eta_0 = 0.0024 \text{ step}^{-1}$ and crystallization timescale $\tau_c = 280$ steps, leading to the update
\begin{equation}
\begin{aligned}
    \mathbf{s}_i(t+1) = \mathrm{clip}_{[-1, 1]^D}\!\Big(& \mathbf{s}_i(t) + \big[\eta_{\mathrm{cryst}, i}(t) + \eta_{\mathrm{sust}, i}(t)\big] \\
    & \times \big(\mathbf{x}_i(t+1) - \mathbf{s}_i(t)\big)\Big).
\end{aligned}
\label{eq:methods_conviction}
\end{equation}

\subsection{Editorial Gatekeeping in Mass Media}
\label{subsec:mass_media}
Media conglomerates $g \in \{1, \dots, M\}$ operate across five modalities $v \in \{0, 1, 2, 3, 4\}$ (Broadcast Television, Radio, Print, Outdoor Signage, and Campaign Rallies) with demographic reach probabilities $P_v(a_i, \mathrm{SES}_i)$. When a factual event $\mathbf{f}_k(t) \in \Omega$ occurs, outlet $g$ projects the notice toward its ideological center $\mathbf{m}_g \in \Omega$ through tendentiousness parameter $\tau_g \in [0, 1]$, according to
\begin{equation}
\begin{aligned}
    \mathbf{y}_{k,g,v}(t) &= (1 - \eta_{g,v})\,\mathbf{f}_k(t) + \eta_{g,v}\,\mathbf{m}_g + \boldsymbol{\xi}_{k,g,v}, \\
    \eta_{g,v} &= \mathrm{clip}_{[0, 1]}(\tau_g \mu_v),
\end{aligned}
\label{eq:methods_media_projection}
\end{equation}
where $\mu_v$ is modal baseline reach and $\sigma_{\mathrm{media}} = 0.03$. Published notices decay exponentially as $I_v(t+1) = I_v(t)\exp(-\delta_v)$. Citizens select active notices via distance-weighted softmax selective exposure,
\begin{equation}
    P(g, v \mid i) = \frac{P_{iv} \, I_{gv}(t) \, e^{-\beta_{\mathrm{m}} d_{igv}(t)}}{\sum_{g',v'} P_{iv'} \, I_{g'v'}(t) \, e^{-\beta_{\mathrm{m}} d_{ig'v'}(t)} + \kappa_0},
    \label{eq:methods_selective_exposure}
\end{equation}
updating opinions within credibility-expanded confidence bounds $d(\mathbf{x}_i, \mathbf{y}_v) \le \epsilon_i (1 + 1.10\,\gamma_v)$ with coupling efficiency $\kappa_{\mathrm{trad}} = 0.08$.

\subsection{Dual-Valence Algorithmic Recommendation}
\label{subsec:digital_platforms}
Digital media encompasses four platform modalities $p \in \{0, 1, 2, 3\}$ (Instagram/TikTok, YouTube/X, Facebook, and WhatsApp). Content creators accumulate followers via preferential attachment $\Pi_k \propto k_i (0.65 + 0.70\tau_{\mathrm{inf}, k})$, which produces scale-free in-degrees $P(k) \sim k^{-1.15}$ with an empirical Gini coefficient of $0.73$. Feed recommendation engines score candidate posts via homophily weighting $q_{ij} \propto (d(\mathbf{x}_i, \mathbf{y}_j) + \delta)^{-\beta_{\mathrm{cur}}}(1 + \mu_{\mathrm{emo}} e_j)\ln(1 + \mathrm{pop}_j)(0.5 + 0.5\ln(1 + k_{\mathrm{inf}}))$ with homophily exponent $\beta_{\mathrm{cur}} = 2.0$ and creator emotionality weight $\mu_{\mathrm{emo}} = 1.5$.

Reactions branch into affirmative likes ($d \le \epsilon_{\mathrm{like}, i}$, assimilating toward $\mathbf{y}_j$) and moral outrage ($d \ge \tau_{\mathrm{dislike}, i}$, triggering backfire repulsion away from $\mathbf{y}_j$). Algorithmic relevance scores compound dual-valence engagement symmetrically as
\begin{equation}
    R_m = \mathrm{Likes}_m + \omega \cdot \mathrm{Dislikes}_m, \qquad \omega = 1.0,
    \label{eq:methods_relevance}
\end{equation}
This symmetric weighting maximizes viral discovery along the diagonal of controversy. Information propagates through a three-stage cascade spanning open-feed broadcast injection, algorithmic recommendation discovery, and dark-social closed micro-group forwarding ($K \le 256$ participants per group, with peer trust multiplier $\tau_{\text{WA}} = 0.85$ and a $2.8\times$ persuasion efficiency advantage).

\subsection{Asymmetric Affective Voting Calculus}
\label{subsec:voting_calculus}
Candidates $k \in \{L, F, O\}$ occupy fixed ideological coordinates $\mathbf{c}_k \in \Omega$. Candidate affinity $A_{ik}(t)$ and aversion $R_{ik}(t)$ evolve with asymmetric parameters ($\lambda_A = 0.12, \delta_A = 0.06; \lambda_R = 0.14, \delta_R = 0.04$),
\begin{align}
    A_{ik}(t+1) &= (1 - \delta_A) A_{ik}(t) \nonumber \\
    &\quad + \lambda_A \Big[\max\!\Big(0, \frac{0.62 - d_{ik}}{0.62}\Big)\Big]^2, \label{eq:methods_adm} \\
    R_{ik}(t+1) &= (1 - \delta_R) R_{ik}(t) \nonumber \\
    &\quad + \lambda_R \Big[\max\!\Big(0, \frac{d_{ik} - 0.78}{1.22}\Big)\Big]^{1.35}. \label{eq:methods_rej}
\end{align}
Net candidate utility is penalized quadratically by aversion,
\begin{equation}
    U_{ik}(t) = A_{ik}(t) - \theta_{\mathrm{rej}} R_{ik}(t)^2, \qquad \theta_{\mathrm{rej}} = 1.65.
    \label{eq:methods_utility}
\end{equation}
Conditional on voting, citizen $i$ chooses candidate $k \in \{L, F, O\}$ via the multinomial logit formulation,
\begin{equation}
    P(\text{Vote } k \mid i, \text{turnout}) = \frac{\exp(\beta_{\mathrm{vote}} U_{ik}(t))}{\sum_{m \in \{L, F, O\}} \exp(\beta_{\mathrm{vote}} U_{im}(t))},
    \label{eq:methods_logit}
\end{equation}
with choice sensitivity $\beta_{\mathrm{vote}} = 3.2$. Turnout probability is modeled via a logistic threshold,
\begin{equation}
\begin{aligned}
    P(\text{Turnout} \mid i) &= \Big[1 + \exp\!\big(-\mu_{\mathrm{turn}}\big[\max_k U_{ik}(t) \\
    &\quad - (\tau_{\mathrm{abst}} - 0.22\,\mathrm{SES}_i + \delta_{\mathrm{transit}})\big]\big)\Big]^{-1}.
\end{aligned}
\label{eq:methods_turnout}
\end{equation}
with steepness $\mu_{\mathrm{turn}} = 4.5$ and free transit subsidy $\delta_{\mathrm{transit}} = -0.15$.

\section{Empirical Setting and Simulation Framework}
\label{sec:empirical_setting}

\subsection{The 2026 Brazilian Presidential Contest}
\label{subsec:brazil_2026}
The 2026 Brazilian presidential contest provides a detailed longitudinal record of electoral competition under divergent communicative regimes (Fig.~\ref{fig:fig1_election_data}). Under Brazilian electoral law (Lei n\textordmasculine~9.504/1997), presidential elections employ a two-round majoritarian system with strict statutory regulations governing campaign finance and broadcast airtime. In the first round, incumbent president Luiz Inácio Lula da Silva assembled a broad seven-party coalition securing 5m~31s per daily broadcast block ($44.1\%$ of total airtime, 433 thirty-second insertions), compared to 4m~20s ($34.7\%$, 339 insertions) for Senator Flávio Bolsonaro; this produced a $+27.3\%$ broadcast endowment advantage for the incumbent (Fig.~\ref{fig:fig1_election_data}d). Financially, while both campaigns adhered to the statutory first-round spending ceiling of $\mathrm{R}\$$~88.9 million (Fig.~\ref{fig:fig1_election_data}e), they diverged sharply in their revenue models (Fig.~\ref{fig:fig1_election_data}f). Specifically, Lula's campaign was $96.3\%$ funded by public transfers from the Special Campaign Financing Fund (FEFC; $\mathrm{R}\$$~123.8M of $\mathrm{R}\$$~128.6M declared), whereas Flávio's campaign deployed a hybrid model raising an unprecedented $\mathrm{R}\$$~43.2M in physical individual donations and micro-crowdfunding, which yielded a $9.0\times$ private capital advantage over Lula ($\mathrm{R}\$$~4.8M).

Throughout the 50-day first-round campaign, voter intentions remained stable (Fig.~\ref{fig:fig1_election_data}b). Longitudinal tracking surveys by Datafolha \citep{datafolha2026r1, datafolha2026r2} placed Lula ahead in a stable band of $41.8\%$--$45.0\%$ of valid votes against $37.5\%$--$41.0\%$ for Flávio, with third-way contenders holding $14.0\%$--$22.0\%$. Concurrently, candidate rejection rates remained mutual and stable between $45\%$ and $47\%$ (Fig.~\ref{fig:fig1_election_data}c). However, the first-round ballot on October 4 produced an abrupt empirical shock (Fig.~\ref{fig:fig1_election_data}a). Official Superior Electoral Court (TSE) ballot tallies \citep{tse2026resultados} recorded Flávio Bolsonaro in first place with $47.03\%$ of valid votes ($55.7$ million votes), surpassing Lula ($45.16\%$, $53.5$ million votes) as third-party support collapsed to $7.81\%$ ($9.3$ million votes). Concurrently, voter abstention and blank ballots surged to a historic high of $24.70\%$ ($20.90\%$ abstention, $4.80\%$ blank/null; Fig.~\ref{fig:fig1_election_data}a). On election night, conservative digital channels mobilized with high intensity, recording $3.90$ million daily interactions compared to Lula's $0.17$ million, and achieved a $22.9\times$ engagement advantage (Fig.~\ref{fig:fig1_election_data}h).

\begin{figure*}[t!]
\centering
\includegraphics[width=\textwidth]{figures/figure0_jan_electiondata.pdf}
\caption{\textbf{Asymmetric campaign endowments and late-stage digital surges drove a two-round electoral reversal in Brazil in 2026.} 
\textbf{a}, Official election returns and voter turnout across Round~1 and Round~2 \citep{tse2026resultados}. 
\textbf{b}, Longitudinal tracking of valid vote intentions across sequential survey waves benchmarked against official ballot counts \citep{datafolha2026r1, datafolha2026r2}. 
\textbf{c}, Candidate rejection trajectories during both electoral rounds \citep{datafolha2026r1, datafolha2026r2}. 
\textbf{d}, Daily statutory broadcast airtime (HGPE) allocated to each candidate across rounds. 
\textbf{e}, Public campaign fund distributions (FEFC) and statutory spending limits. 
\textbf{f}, Revenue compositions contrasting public party transfers with private micro-donations \citep{tse2026divulgacand}. 
\textbf{g}, Candidate digital popularity dynamics measured via the Quaest Digital Popularity Index \citep{quaest2026ipd}. 
\textbf{h}, Daily social media interaction volume across candidate networks during major campaign milestones \citep{bites2026digital}. 
\textbf{i}, Cross-platform digital audience and surrogate bench scale across major open networks and private messaging channels.}
\label{fig:fig1_election_data}
\end{figure*}

\subsection{Hierarchical Empirical Calibration}
\label{subsec:calibration}
To maintain methodological transparency, all model parameters are partitioned into three distinct classes detailed in Appendix A (Tables~\ref{tab:supp_model_parameters} and \ref{tab:supp_media_parameters}).

First, fixed institutional and demographic parameters are grounded directly in empirical benchmark sources without free optimization. Citizen ages $a_i \in [16, 90]$ are sampled from the official IBGE 2022 Census pyramid with Gaussian smoothing ($\sigma_a = 1.35$~years), socioeconomic status follows $\mathrm{SES}_i \sim \mathrm{Beta}(2.0, 3.0)$, and average household size is set to $2.79$ residents per domicile in Layer $\mathcal{L}_1$ \citep{ibge2022censo}. Workplace and educational cohort sizes ($\langle k^{[2]} \rangle = 10$) reflect registries from the INEP School Census and MTE/RAIS labor databases \citep{inep2023censo, mte2023rais}. Traditional media reach profiles across television, radio, print, and outdoor channels are drawn from Kantar IBOPE audience surveys \citep{kantar2023inside}, while digital platform penetration metrics across Instagram, TikTok, YouTube, X, Facebook, and WhatsApp ($K \le 256$ forwarding limit) follow the Reuters Institute Digital News Report \citep{reuters2023digital}. Statutory electoral parameters adhere to Superior Electoral Court regulations under Lei n\textordmasculine~9.504/1997, specifying daily HGPE broadcast airtime (Lula 5m~31s versus Fl\'avio 4m~20s in Round~1, and strict 5m~00s parity in Round~2), presidential campaign expenditure ceilings, and public party transfers from the FEFC fund. Candidate ideological coordinates in concept space $\Omega$ place Lula at $\mathbf{c}_L = (-0.25, -0.20, -0.35, -0.15)^{\top}$, Fl\'avio Bolsonaro at $\mathbf{c}_F = (+0.65, +0.70, +0.60, +0.65)^{\top}$, and Third-Way centrists at $\mathbf{c}_O = (+0.20, +0.25, +0.10, +0.30)^{\top}$, derived from party manifesto analyses and legislative roll-call records.

Second, behavioral parameters are anchored in established cognitive psychology and complex network benchmarks. Social influence weights and the row-stochastic cognitive issue coupling matrix $\mathbf{C}$ follow the Friedkin--Johnsen framework \citep{friedkin1990social, friedkin2016network}. Bounded-confidence thresholds ($\bar{\epsilon} = 0.45$) and relative-agreement uncertainty overlap follow standard formulations \citep{deffuant2000mixing, deffuant2002how}, complemented by biased assimilation weighting \citep{lord1979biased, dandekar2013biased}. Loss-aversion framing adopts Kahneman--Tversky prospect theory multipliers with $\lambda_{\mathrm{loss}} \in [2.0, 2.5]$ \citep{kahneman1979prospect}. Network parameters incorporate Watts--Strogatz small-world rewiring ($p_{\mathrm{ws}} = 0.15$) in friendship ties \citep{watts1998collective}, Granovetter weak ties in transit encounters \citep{granovetter1973strength}, and Barab\'asi--Albert preferential attachment scaling ($\gamma_{\mathrm{pl}} = 1.15$, Gini coefficient $0.73$) for digital influencers \citep{barabasi1999emergence, bianconi2001competition}. The elastic tethering function $g_i^{\mathrm{eff}}$ implements Sherif--Hovland latitudes of rejection around core convictions \citep{sherif1961social, howe2017attitude}.

Third, free parameters fitted to longitudinal tracking surveys and official election results are estimated via bounded grid search and Markov Chain Monte Carlo sampling. These comprise the discrete integration step size ($\Delta t = 0.15$) and cognitive noise ($\sigma_{\xi} = 0.012$), baseline conviction adaptation drift ($\eta_0 = 0.0024 \text{ step}^{-1}$) and crystallization horizon ($\tau_c = 280$ simulation steps), and the factual pull coefficient ($\gamma_0 = 0.055 \text{ step}^{-1}$). They also encompass candidate admiration and aversion kinetic rates ($\lambda_A = 0.12, \delta_A = 0.06; \lambda_R = 0.14, \delta_R = 0.04$), which establish the kinetic asymmetry ratio $\lambda_R / \delta_R = 3.5\times$ and aversion penalty $\theta_{\mathrm{rej}} = 1.65$. Finally, voting choice parameters specify logit sensitivity $\beta_{\mathrm{vote}} = 3.2$, base abstention threshold $\tau_{\mathrm{abst}} = 0.38$, turnout logistic steepness $\mu_{\mathrm{turn}} = 4.5$, and transit subsidy impact $\delta_{\mathrm{transit}} = -0.15$.

The optimization objective minimizes weighted squared error across all survey waves and official election totals,
\begin{equation}
\begin{aligned}
    L(\boldsymbol{\theta}) = \; & \sum_{w=1}^W \sum_{k \in \{L, F, O\}} \big(V_k^{\mathrm{model}}(t_w; \boldsymbol{\theta}) - V_k^{\mathrm{Datafolha}}(t_w)\big)^2 \\
    & + \omega_{\mathrm{ballot}} \sum_{k} \big(V_k^{\mathrm{model}}(t_{\mathrm{election}}; \boldsymbol{\theta}) - V_k^{\mathrm{TSE}}\big)^2,
\end{aligned}
\label{eq:methods_loss_function}
\end{equation}
with ballot weight $\omega_{\mathrm{ballot}} = 3.0$. This parameterization achieves $\mathrm{RMSE} = 0.85$ percentage points ($R^2 = 0.9955$) across Round~1 and $\mathrm{RMSE} = 0.56$ percentage points ($R^2 = 0.9128$) across Round~2.

\subsection{High-Performance Simulation Architecture}
\label{subsec:hpc}
Simulations were executed across $S = 100$ independent Monte Carlo realizations at population scale $N = 100{,}000$ agents using just-in-time (JIT) compiled Numba 0.59 kernels in Python 3.11 with parallel multi-threading across 64 CPU cores (AMD EPYC 7763, 2.45 GHz, 512 GB RAM). All agent attributes are laid out as contiguous 1D and 2D C-ordered NumPy arrays in CPU cache memory, with network topology represented in CSR format. At $N = 100{,}000$, a 100-step trajectory executes in $4.2$ seconds. Municipal scaling experiments ($N = 10^4$ to $10^7$) confirm that finite-size stochasticity vanishes into a scale-invariant thermodynamic limit for $N \ge 10^5$.

Macroscopic system trajectories are evaluated using four statistical metrics. Polarization is monitored through the multidimensional Esteban--Ray index \citep{esteban1994measurement},
\begin{equation}
    P_{\mathrm{ER}}(t) = \kappa_{\mathrm{ER}} \sum_{c, c'} \pi_c^{1 + \alpha_{\mathrm{ER}}} \pi_{c'} d(\bar{\mathbf{x}}_c, \bar{\mathbf{x}}_{c'}),
    \label{eq:er_index}
\end{equation}
with sensitivity $\alpha_{\mathrm{ER}} = 1.30$ and scaling $\kappa_{\mathrm{ER}} = 1.0$. Distributional modality is evaluated via Sarle's Bimodality Coefficient $\mathrm{BC} = (\gamma^2 + 1) / [\kappa + 3 (n-1)^2 / ((n-2)(n-3))]$. Effective ideological dimensionality is tracked by the Participation Ratio of the belief covariance matrix, $\mathrm{PR}(\Sigma) = (\mathrm{Tr}\,\Sigma)^2 / \mathrm{Tr}(\Sigma^2) \in [1, D]$ \citep{roy1953useful}. Finally, issue constraint is quantified by the mean absolute pairwise correlation $\langle |\rho| \rangle = \frac{2}{D(D-1)} \sum_{k < l} |\rho_{kl}|$. All time-series trajectories report ensemble means with $95\%$ confidence intervals ($\bar{\mu} \pm 1.96\,\hat{\sigma}/\sqrt{S}$) or full simulation ranges $[\min, \max]$.

\section{Experimental Results}
\label{sec:results}

\subsection{Scale-Invariant Multiplex Relaxation}
\label{subsec:results_baseline}
To establish the baseline stability of public attitudes prior to media exposure or algorithmic amplification, I first quantify the relaxation physics of the urban interpersonal network (Fig.~\ref{fig:fig2_interpersonal_network}). Calibrated to IBGE 2022 National Population Census demography \citep{ibge2022censo}, the model situates citizens in a four-dimensional concept hypercube $\Omega = [-1, 1]^4$ \citep{converse1964nature}. Individual behavioral parameters, cognitive plasticity $\alpha_i$, bounded confidence threshold $\epsilon_i$, attitude uncertainty $u_i$, and confirmation bias $b_i$, are conditioned on empirical age and socioeconomic status (SES; Fig.~\ref{fig:fig2_interpersonal_network}a--c). In agreement with the aging-stability hypothesis \citep{alwin1991aging, kiley2020measuring}, cognitive plasticity $\alpha(a)$ undergoes monotonic decay with age (Fig.~\ref{fig:fig2_interpersonal_network}a), while older and lower-SES cohorts exhibit narrower confidence bounds $\epsilon_i$ and elevated confirmation bias $b_i$.

\begin{figure*}[t!]
\centering
\includegraphics[width=\textwidth]{figures/figure1_jan_nonmedianetwork.pdf}
\caption{\textbf{Decentralized deliberation across urban multiplex networks converges toward a scale-invariant centrist equilibrium while buffering core convictions against acute shocks.} 
\textbf{a}, Empirical demographic age distribution and age-dependent cognitive plasticity $\alpha(a)$. 
\textbf{b}, Socioeconomic status distribution fitted to census benchmarks. 
\textbf{c}, Cohort distributions of individual cognitive and behavioral parameters. 
\textbf{d}, Layered multiplex topology spanning households, workplaces, congregations, friendship circles, and transit ties. 
\textbf{e}, Inter-cohort mixing matrix across adult age brackets. 
\textbf{f}, Structural homophily and ideological assortativity across network layers. 
\textbf{g}, Two-dimensional projection of initial versus converged convictions alongside exogenous factual anchors. 
\textbf{h}, Cross-issue alignment trajectories across urban population scales. 
\textbf{i}, Asymptotic issue alignment and effective belief dimensionality versus city population scale. 
\textbf{j}, Stationary marginal probability densities across ideological dimensions. 
\textbf{l}, Dynamical relaxation of expressed opinions and internalized convictions following an exogenous shock. 
\textbf{m}, Age-stratified opinion--conviction tether displacement across time. 
\textbf{n}, Geometric snapshot of attitude space at peak shock excursion comparing displaced opinions with anchored convictions.}
\label{fig:fig2_interpersonal_network}
\end{figure*}

Interpersonal interactions unfold over a five-layer multiplex topology $\mathcal{M} = \{\mathcal{L}_1, \dots, \mathcal{L}_5\}$ \citep{boccaletti2014structure, kivela2014multilayer} capturing household cohabitation ($\mathcal{L}_1$, mean size $2.79$, weight $w_1 = 0.36$), workplace/school cohorts ($\mathcal{L}_2$, $\langle k^{[2]} \rangle = 10$, $w_2 = 0.26$), religious and community congregations ($\mathcal{L}_3$, $\langle k^{[3]} \rangle = 8$, $w_3 = 0.14$), small-world friendship circles ($\mathcal{L}_4$, $\langle k^{[4]} \rangle = 6$, $w_4 = 0.14$), and casual transit ties ($\mathcal{L}_5$, $\langle k^{[5]} \rangle = 4$, $w_5 = 0.10$) (Fig.~\ref{fig:fig2_interpersonal_network}d). The empirical age contact matrix $M_{ab}$ balances generational stratification in workplaces against cross-cohort mixing in households and religious institutions (Fig.~\ref{fig:fig2_interpersonal_network}e). This architecture generates a structural trade-off between age homophily $h_{\mathrm{age}}$ and ideological assortativity $r_{\mathrm{ideol}}$ (Fig.~\ref{fig:fig2_interpersonal_network}f). While friendship ($\mathcal{L}_4$, $r_{\mathrm{ideol}} = 0.38$) and religious groups ($\mathcal{L}_3$, $r_{\mathrm{ideol}} = 0.46$) promote ideological sorting, transit encounters ($\mathcal{L}_5$, $r_{\mathrm{ideol}} \approx 0.02$) and domestic ties ($\mathcal{L}_1$, $r_{\mathrm{ideol}} = 0.24$) act as cross-cutting bridges that inhibit macroscopic topological fragmentation.

In the absence of external media distortion, decentralized deliberation over this multiplex substrate systematically dampens ideological extremity (Fig.~\ref{fig:fig2_interpersonal_network}g). Although initial convictions $\mathbf{s}_i(0)$ are widely dispersed across $[-1, 1]^4$, bounded-confidence exchanges pull terminal convictions $\mathbf{s}_i(T)$ toward a cohesive centrist consensus around the origin $\mathbf{0}$ \citep{castellano2009statistical}. Scaling analysis across four orders of magnitude in urban population ($N = 10^4$ to $10^7$) indicates that this macroscopic convergence is scale-invariant (Fig.~\ref{fig:fig2_interpersonal_network}h--j). For populations $N \ge 10^5$, finite-size stochasticity vanishes into a thermodynamic limit where terminal issue correlation ($\max |\rho| \approx 0.08, \langle |\rho| \rangle \approx 0.03$) and effective belief dimensionality ($\mathrm{PR}(\Sigma) \approx 3.75$ out of $D=4$) stabilize at identical asymptotic values (Fig.~\ref{fig:fig2_interpersonal_network}i), while stationary marginal densities $p(s_d)$ remain unimodal and Gaussian-like across all dimensions (Fig.~\ref{fig:fig2_interpersonal_network}j).

This formulation resolves the dynamical coupling between fast expressed opinions $\mathbf{x}_i(t)$ and slow internalized convictions $\mathbf{s}_i(t)$ \citep{zaller1992nature, howe2017attitude} (Fig.~\ref{fig:fig2_interpersonal_network}l--n). Under an acute exogenous factual shock, mean opinion $\bar{\mathbf{x}}(t)$ responds rapidly ($\tau_x \approx 0.15$), whereas mean conviction $\bar{\mathbf{s}}(t)$ exhibits strong structural inertia ($\tau_s / \tau_x \approx 11.3$), shifting only marginally before pulling opinions back toward baseline during post-shock relaxation (Fig.~\ref{fig:fig2_interpersonal_network}l). This separation of timescales induces an elastic tether displacement $\|\mathbf{x}_i(t) - \mathbf{s}_i(t)\|$ that varies across demographics (Fig.~\ref{fig:fig2_interpersonal_network}m). Specifically, younger citizens experience substantial opinion displacement due to higher cognitive plasticity, while older citizens remain tightly tethered to crystallized convictions \citep{alwin1991aging}. A cross-sectional snapshot of the belief hypercube at peak shock excursion ($t_{\mathrm{shock}}$; Fig.~\ref{fig:fig2_interpersonal_network}n) illustrates this instantaneous decoupling between displaced opinions and stationary convictions. This separation explains mechanistically how acute campaign events can trigger large short-term shifts in voting intention without destabilizing core ideological convictions.

\subsection{Media Market Concentration and Polarization}
\label{subsec:results_media}
Introducing the second communicative substrate, institutional mass media channels structure collective public attention via centralized broadcast scheduling and editorial gatekeeping \citep{mccombs1972agenda} (Fig.~\ref{fig:fig3_traditional_media}). I parameterize five distinct modalities, consisting of Broadcast Television ($v = 0$), Radio ($v = 1$), Print and Legacy News Portals ($v = 2$), Outdoor Signage ($v = 3$), and Campaign Rallies ($v = 4$) (Fig.~\ref{fig:fig3_traditional_media}a). Modalities differ in their characteristic attention decay rate $\delta_v$ and institutional credibility $\gamma_v$. Written Press combines high institutional credibility ($\gamma = 0.82$) with slow editorial decay ($\delta = 0.18$), Broadcast Television provides mass reach with intermediate persistence ($\gamma = 0.75, \delta = 0.30$), and Campaign Rallies deliver high emotional mobilization pulses ($\gamma = 0.80$) subject to rapid post-event decay ($\delta = 0.70$). Audience penetration curves $P_v(\mathrm{age})$ calibrated to Kantar IBOPE 2023 surveys \citep{kantar2023inside} indicate that linear television viewing increases monotonically with age, dominating news consumption among senior citizens, whereas Radio and Outdoor signage reach daily commuters (Fig.~\ref{fig:fig3_traditional_media}b).

\begin{figure*}[t!]
\centering
\includegraphics[width=\textwidth]{figures/figure2_jan_medianetwork.pdf}
\caption{\textbf{Editorial tendentiousness and market concentration nonlinearly govern collective factual accuracy and societal polarization.} 
\textbf{a}, Decay rates and credibility parameters across broadcast and print modalities. 
\textbf{b}, Demographic reach curves across age brackets. 
\textbf{c}, Editorial projection of factual events toward conglomerate ideological anchors under objective versus tendentious coverage. 
\textbf{d}, Public opinion trajectory following a factual shock under varied media market configurations. 
\textbf{e}, Macroscopic attitude displacement and Esteban--Ray polarization across market concentration regimes. 
\textbf{f}, Operational timescale separation between transient media pulses and persistent interpersonal network deliberation.}
\label{fig:fig3_traditional_media}
\end{figure*}

Media conglomerates project factual notices toward editorial ideological centers $\mathbf{m}_g \in \Omega$ through their tendentiousness parameter $\tau_g \in [0, 1]$ \citep{mullainathan2005market, gentzkow2006media}. In attitude space (Fig.~\ref{fig:fig3_traditional_media}c), low tendentiousness ($\tau = 0.10$) confines published notices to a tight envelope around true factual coordinates $\mathbf{f}_k$. Conversely, under high tendentiousness ($\tau = 0.90$), editorial gatekeeping displaces published coverage toward the conglomerate's ideological anchor $\mathbf{m}_g$ \citep{puglisi2015newspaper}, substituting factual reality with partisan narratives.

Market structure nonlinearly governs collective responsiveness following an acute factual shock ($\mathbf{f}^* = +0.80$ introduced at $t \ge 25$; Fig.~\ref{fig:fig3_traditional_media}d). Without broadcast media ($M=0$), information diffuses sluggishly through word-of-mouth alone, producing negligible collective adjustment ($\bar{x}_1 \approx 0.00 \to 0.03$). Factual broadcasting ($\tau = 0.10$) rapidly synchronizes public opinion toward empirical reality ($\bar{x}_1 \to +0.50$). Conversely, a biased media monopoly ($M=1, \tau = 0.90$ anchored at $\mathbf{m}_1 = -0.60$) actively drives public opinion in the opposite direction ($\bar{x}_1 \to -0.24$), inducing larger systemic error than the absence of media. Across market configurations ($M \in \{0, 1, 2, 3, 4, \text{Many}\}$; Fig.~\ref{fig:fig3_traditional_media}e), increasing outlet plurality dilutes monopoly capture (Fig.~\ref{fig:fig3_traditional_media}e, top). However, societal division, measured via the Esteban--Ray polarization index $P_{\mathrm{ER}}$ \citep{esteban1994measurement}, peaks sharply under a tendentious duopoly ($M=2$; Fig.~\ref{fig:fig3_traditional_media}e, bottom), which bifurcates citizens into polarized information silos. Only pluralistic markets ($M \ge 4$) disperse partisan sorting and restore low societal polarization comparable to the baseline urban network.

Coupling the broadcast layer to the interpersonal network reveals a fundamental separation of operational timescales (Fig.~\ref{fig:fig3_traditional_media}f). Broadcast notices display high instantaneous reach but rapid temporal decay ($\tau_{\mathrm{media}} \approx 4.8$~steps) \citep{downs1972up}. In contrast, interpersonal deliberation across the multiplex network operates on a substantially slower relaxation timescale ($\tau_{\mathrm{network}} \approx 19.0$~steps) \citep{friedkin1990social}. Consequently, when broadcast channels introduce a salient campaign event, its societal impact is sustained long after the initial broadcast window closes, as citizens continue to process the information across household and peer ties.

\subsection{Algorithmic Outrage Amplification}
\label{subsec:results_digital}
Completing the communication architecture, the digital layer introduces algorithmic social platforms and closed encrypted messaging networks \citep{cinelli2021echo, tucker2018social} (Fig.~\ref{fig:fig4_social_media}). Calibrated to national surveys \citep{reuters2023digital}, age access profiles $P_p(\mathrm{age})$ define four digital modalities (Fig.~\ref{fig:fig4_social_media}a), spanning short-form visual discovery feeds (Platform~0, Instagram and TikTok) that reach over $88\%$ of youth cohorts; video commentary and microblogging (Platform~1, YouTube and X) engaging politically attentive citizens; reciprocal social feeds (Platform~2, Facebook) concentrating among mature adults; and closed encrypted messaging (Platform~3, WhatsApp and Telegram) operating as a universal communications backbone ($>80\%$ penetration across all adult brackets) structured into private micro-groups ($K \le 256$~members).

\begin{figure*}[t!]
\centering
\includegraphics[width=0.88\textwidth]{figures/figure3_jan_socialmedianetwork.pdf}
\caption{\textbf{Engagement-maximizing recommendation algorithms amplify polarizing content and sustain multi-stage cascades into encrypted messaging networks.} 
\textbf{a}, Demographic access probabilities across digital platform modalities and broadcast television. 
\textbf{b}, Scale-free in-degree follower distribution and structural hierarchy among digital influencers. 
\textbf{c}, Algorithmic relevance landscape across dual-valence engagement dimensions, showing peak amplification along the diagonal of controversy. 
\textbf{d}, Multi-stage diffusion lifecycle linking open discovery feeds, public commentary platforms, and encrypted messaging groups. 
\textbf{e}, Hybrid agenda-setting feedback loop contrasting legacy editorial turnaround with high-velocity algorithmic platform diffusion. 
\textbf{f}, Long-term societal polarization trajectories under varying degrees of content creator tendentiousness.}
\label{fig:fig4_social_media}
\end{figure*}

Within open digital feeds, follower accumulation produces extreme structural inequality (Fig.~\ref{fig:fig4_social_media}b). Preferential attachment yields a scale-free in-degree distribution $P(k) \sim k^{-\gamma_{\mathrm{pl}}}$ with exponent $\gamma_{\mathrm{pl}} \approx 1.15$ and an empirical Gini coefficient of $0.73$ (Fig.~\ref{fig:fig4_social_media}b). Under this topology, the top $5\%$ of elite creators capture $54.5\%$ of direct audience reach, whereas the bottom $50\%$ command only $8.4\%$. This hyper-concentration converts decentralized networks into a star-like broadcasting system; when an elite influencer publishes a narrative, their message bypasses local interpersonal clustering to reach hundreds of thousands of individual algorithmic feeds simultaneously.

Feed recommendation algorithms rank content by optimizing for aggregate user engagement rather than factual accuracy \citep{vosoughi2018spread, crockett2017moral}. Computing relevance scores symmetrically across positive and negative interactions creates a structural amplification along the diagonal of controversy (Fig.~\ref{fig:fig4_social_media}c). Posts provoking both in-group admiration and out-group outrage receive maximal algorithmic distribution, rewarding polarizing moral-emotional framing \citep{brady2017emotion, crockett2017moral}.

This architecture generates a three-stage cross-platform cascade lifecycle (Fig.~\ref{fig:fig4_social_media}d). Diffusion initiates in Stage~1 on visual discovery feeds (Instagram/TikTok), surfacing novel memes to high-plasticity youth within $10$~hours. In Stage~2, narratives migrate to commentary platforms (YouTube/X), where quote-dunking triggers secondary viral waves between hours $15$ and $25$. In Stage~3, content penetrates closed messaging groups (WhatsApp/Telegram). Forwarded across trusted peer contacts ($\tau_{\text{WA}} = 0.85$, $2.8\times$ persuasion efficiency), dark-social cascades bypass platform moderation and persist long after open-feed attention has decayed. Coupling the digital layer to traditional newsrooms creates a hybrid agenda-setting feedback loop \citep{chadwick2017hybrid, harder2017intermedia} (Fig.~\ref{fig:fig4_social_media}e). While investigative journalism requires an editorial turnaround of $\tau_{\mathrm{turnaround}} \approx 1$--$2$~days, digital platforms respond within hours ($\tau_{\mathrm{social}} \approx 0.1$--$0.3$~days); this rapid response generates social volume that compels broadcast media to maintain follow-up coverage. When influencers exploit this dynamic with high tendentiousness ($\tau_{\mathrm{inf}} = 0.90$), societal polarization climbs to an entrenched plateau of $P_{\mathrm{ER}} = 0.18$ (Fig.~\ref{fig:fig4_social_media}f).

\subsection{First-Round Reversal via Dark-Social Blitzes}
\label{subsec:results_first_round}
The convergence of multiplex interpersonal clustering, traditional media gatekeeping, and algorithmic outrage explains the dynamics of the first-round election (Fig.~\ref{fig:fig5_first_round}). Calibrating initial ideological identification to national survey benchmarks \citep{datafolha2024identificacao} establishes four distinct voting blocs (Fig.~\ref{fig:fig5_first_round}a), comprising Lulistas ($\sim 40\%$, comprising $28\%$ convinced identifiers and $12\%$ non-convinced sympathizers), Bolsonaristas ($\sim 34\%$, $23\%$ convinced and $11\%$ sympathizers), Centrists ($\sim 19\%$), and unaligned non-voters ($\sim 7\%$). In concept space $\Omega$, Lula is positioned at Center-Left ($\mathbf{c}_L = [-0.25, -0.20, -0.35, -0.15]^\top$), Flávio Bolsonaro at Far-Right\linebreak ($\mathbf{c}_F = [0.65, 0.70, 0.60, 0.65]^\top$), and Third-Way contenders at Center-Right\linebreak ($\mathbf{c}_O = [0.20, 0.25, 0.10, 0.30]^\top$). Brazilian centrists are culturally conservative on moral and security axes ($d=1, 3$) yet supportive of state welfare and minimum wages ($d=0, 2$), making them a cross-pressured swing constituency \citep{almeida2019evangelicals}.

\begin{figure*}[t!]
\centering
\includegraphics[width=0.82\textwidth]{figures/figure4_jan_election_first.pdf}
\caption{\textbf{Cognitive negativity asymmetry and a 48-hour dark-social blitz overturned an established polling lead in the first round.} 
\textbf{a}, Electorate segmentation into ideological identification blocs calibrated to national survey data. 
\textbf{b}, Longitudinal 50-day campaign trajectories comparing model projections with tracking surveys (diamonds, mean $\pm 1\,\mathrm{SD}$) and official ballot tallies (circles), marking the debate shock ($*$) and pre-election blitz window. 
\textbf{c}, Population-mean candidate admiration and rejection dynamics illustrating the late-campaign aversion surge. 
\textbf{d}, Factorial attribution of vote share shifts across campaign mechanisms, with error bars denoting the full simulation range $[\min, \max]$ across 100 independent Monte Carlo seeds ($N=100{,}000$). 
\textbf{e}, Parametric dose-response curves showing vote share sensitivity to mechanism intensity, with calibrated empirical operating points and benchmark contours. 
\textbf{f}, Stepwise cumulative model addition across successive mechanisms, showing valid vote shares (error bars: full range $[\min, \max]$ across 100 seeds) and lead inversion during the final dark-social blitz.}
\label{fig:fig5_first_round}
\end{figure*}

Simulating the 50-day first round across 100 independent Monte Carlo realizations ($N = 100{,}000$ citizens) reproduces Datafolha survey waves and the official TSE outcome ($47.0\%$ Flávio, $45.2\%$ Lula, $7.8\%$ Others) with high statistical fidelity ($\mathrm{RMSE} = 0.85$ percentage points, $R^2 = 0.9955$; Fig.~\ref{fig:fig5_first_round}b). Under the neutral organic baseline (dashed curves), Lula preserves an uninterrupted lead throughout the campaign, concluding at $48.2\%$ ($57.1\mathrm{M}$ votes) versus $38.8\%$ for Flávio ($46.0\mathrm{M}$ votes). In the fitted model (solid curves), Lula leads securely through day 47 ($48.2\%$ vs $38.8\%$). During the final 48-hour sprint (shaded amber band), a synchronized multi-platform offensive produces a sharp vote shift of $+8.20$ percentage points ($+9.72\mathrm{M}$ votes) that propelled Flávio ahead on election day.

The behavioral mechanism driving this late-stage divergence is a cognitive negativity asymmetry governing candidate admiration and rejection (Fig.~\ref{fig:fig5_first_round}c). While admiration accumulates slowly and decays rapidly, candidate aversion accumulates rapidly and decays with high hysteresis. This kinetic asymmetry, where aversion accumulates $3.5$ times faster than it decays, implies that negative shocks inflict persistent psychological scars that resist correction. During the final 48 hours, targeted moral panics elevate Lula's rejection from $0.38$ to $0.50$ (Fig.~\ref{fig:fig5_first_round}c). Because voter utility penalizes candidate aversion more heavily than it rewards positive affinity, this rejection spike depresses swing-voter turnout and shifts cross-pressured centrists toward Flávio.

Factorial attribution isolates the specific contribution of each mechanism toward Flávio's $+8.20\%$ ($+9.72\mathrm{M}$ votes) surge (Fig.~\ref{fig:fig5_first_round}d). Religious Priming (M1, $+1.77\%$ or $+2.10\mathrm{M}$ votes) anchors Evangelical voters in layer $\mathcal{L}_3$, Traditional Media Slant (M2, $+1.20\%$ or $+1.42\mathrm{M}$ votes) reinforces fiscal conservatism, Bot Astroturfing (M3, $+1.35\%$ or $+1.60\mathrm{M}$ votes) manufactures artificial outrage on open feeds, Asymmetric Rejection News (M4, $+1.55\%$ or $+1.84\mathrm{M}$ votes) exploits cognitive negativity bias, and the 48-Hour WhatsApp Blitz (M5, $+2.20\%$ or $+2.61\mathrm{M}$ votes) saturates dark-social groups beyond journalistic verification horizons. Parametric dose-response curves reveal that the WhatsApp blitz exhibits the steepest marginal electoral return ($\mathrm{d}V_F/\mathrm{d}\theta \approx 0.034\%$ per percentage point; Fig.~\ref{fig:fig5_first_round}e). A cumulative build-up test reveals a threshold principle (Fig.~\ref{fig:fig5_first_round}f). Lula preserves his polling lead through stages 1 to 4 (+Church, +Media, +Bots, +Rejection); decisive lead reversal occurs exclusively when the acute dark-social blitz is added in Stage 5, placing Flávio ahead at $47.0\%$ versus $45.2\%$.

\subsection{Graph-Theoretic Formalization of Counter-Tactics}
\label{subsec:janonismo_tactics}
The first-round diagnosis informed the design of the progressive counter-offensive in the second-round runoff (October 4--25, 2026; Fig.~\ref{fig:fig6_second_round}), which operationalized the strategic doctrine of \emph{Janonismo Cultural} \citep{janones2023janonismo}. In counterfactual simulations, if the progressive campaign had maintained its passive first-round posture relying on broadcast airtime and conventional fact-checking, the same-strategy counterfactual (dashed curves, Fig.~\ref{fig:fig6_second_round}a) indicates that Flávio would have consolidated centrist support to win decisively with $53.6\%$ ($63.9\mathrm{M}$ votes) against $46.4\%$ for Lula ($55.4\mathrm{M}$ votes).

\begin{figure*}[t!]
\centering
\includegraphics[width=0.82\textwidth]{figures/figure5_jan_election_second.pdf}
\caption{\textbf{Coordinated socio-technical counter-tactics inverted candidate rejection trajectories and crossed the runoff majority threshold.} 
\textbf{a}, Longitudinal 21-day runoff trajectories contrasting the unadjusted counterfactual against the fitted counter-offensive (envelope: mean $\pm 1\,\mathrm{SD}$ across 100 seeds), benchmarked against tracking polls and official runoff returns. 
\textbf{b}, Candidate affective trajectories showing inversion of rejection and consolidation of admiration across the runoff campaign. 
\textbf{c}, Factorial attribution of vote share turnaround across the six counter-tactics, with error bars denoting the full simulation range $[\min, \max]$ across 100 independent Monte Carlo seeds ($N=100{,}000$). 
\textbf{d}, Multi-tactic parametric dose-response curves and iso-vote contours mapping tactic intensity to vote share gain, showing calibrated empirical operating points. 
\textbf{e}, Cumulative stepwise addition of counter-tactics illustrating progressive vote share gains (error bars: full range $[\min, \max]$ across 100 seeds) and crossing of the majority threshold.}
\label{fig:fig6_second_round}
\end{figure*}

Under the active counter-tactical suite, the campaign deployed six structured interventions formalizing the operational principles of \emph{Janonismo Cultural} into graph-theoretic mechanisms, summarized in Table~\ref{tab:supp_janones_tactics}. 

The first set of interventions operated directly within open algorithmic discovery feeds. Through an outrage-starvation embargo (Tactic T1, Algorithmic Asphyxiation, $+1.14\%$ or $+1.36\mathrm{M}$ votes), allied accounts enforced complete non-engagement toward opponent provocations; this withheld the controversy signals curation algorithms require and reduced opponent feed injection by $48\%$. In parallel, monothematic bandwidth concentration (Tactic T2, Strict Monothematism, $+1.35\%$ or $+1.61\mathrm{M}$ votes) coordinated communication around a single negative theme per 24-hour cycle, leveraging superlinear returns in platform recommendation functions to quadruple viral reach relative to fragmented messaging. The campaign complemented this focus with material loss-aversion framing (Tactic T3, Material Loss Frame, $+1.42\%$ or $+1.69\mathrm{M}$ votes), replacing abstract defenses of democratic norms with concrete economic risks such as welfare rollbacks and wage freezes; this engaged Kahneman--Tversky loss weighting (a 2.25-fold aversion to losses relative to gains) and converted $52.4\%$ of undecided working-class swing voters.

The second set of interventions targeted dark-social distribution channels and physical mobilization. Through molecular encrypted cascades (Tactic T4, WhatsApp Synchronization, $+1.05\%$ or $+1.25\mathrm{M}$ votes), a registered network of $10.0\mathrm{M}$ volunteers dispatched curated multimedia assets into private family and neighborhood groups, using direct peer trust to achieve a 2.8-fold persuasion multiplier over broadcast posts. Concurrently, physical turnout locking (Tactic T5, Mass Capital Rallies, $+0.88\%$ or $+1.05\mathrm{M}$ votes) organized large assemblies across 27 state capitals ($1.80\mathrm{M}$ attendees) and combined with municipal transit subsidies to eliminate voting travel friction, cementing partisan turnout. Finally, attention-ceiling saturation (Tactic T6, 24-Hour Continuous Livestream, $+0.80\%$ or $+0.95\mathrm{M}$ votes) sustained continuous broadcasts throughout the final 48 hours, absorbing $65\%$ of online political attention and suppressing rival late-campaign discovery below $20\%$.

\subsection{Runoff Reversal and the Majority Threshold}
\label{subsec:results_runoff}
Together, these interventions produced a complete inversion of candidate affective evaluations (Fig.~\ref{fig:fig6_second_round}b). Flávio's rejection rose from $\bar{R}_F = 0.40$ to $0.53$, while Lula's rejection dropped from $\bar{R}_L = 0.50$ to $0.41$, alongside a concurrent rise in Lula admiration ($\bar{A}_L = 0.43 \to 0.52$). The calibrated model reproduces the 21-day runoff tracking polls and official election results with high precision ($\mathrm{RMSE} = 0.56$ percentage points, $R^2 = 0.9128$; Fig.~\ref{fig:fig6_second_round}a). Parametric dose-response curves illustrate tactical fungibility across the intervention suite (Fig.~\ref{fig:fig6_second_round}d). A cumulative addition test reveals a decisive majority crossover (Fig.~\ref{fig:fig6_second_round}e). Starting from the Round 1 baseline ($45.16\%$), tactics T1 ($46.30\%$), T2 ($47.65\%$), and T3 ($49.07\%$) steadily close the gap; the $50\%$ majority threshold is crossed upon adding Tactic T4 (reaching $50.12\%$), after which T5 ($51.00\%$) and T6 ($51.80\%$) consolidate the final victory.

\section{Discussion}
\label{sec:discussion}

\subsection{Mechanistic Resolution of Late-Campaign Polling Misses}
The findings presented here provide a generative framework to address a fundamental paradox in contemporary democratic politics, namely, how modern elections can experience sudden, late-stage reversals that diverge from eve-of-election survey forecasts despite the documented stability of voters' foundational political convictions. By developing an empirical digital twin of the 2026 Brazilian presidential contest, the model accounts for how late-campaign volatility can emerge not from the wholesale restructuring of underlying belief systems, but from an acute dynamical decoupling between fast-moving expressed opinions $\mathbf{x}_i(t)$ and slow-moving ideological convictions $\mathbf{s}_i(t)$.

This formulation provides a mechanistic explanation for the two-round polling anomaly observed in Brazil in 2026. While pre-election survey discrepancies in one-round races typically stem from a single late-breaking mobilization, sequential late-stage swings in opposing directions across rounds reflect the dynamic properties of two-timescale opinion systems. Traditional polling methodologies sample top-of-the-head public opinions $\mathbf{x}_i(t)$ at discrete intervals under an implicit assumption of quasi-static equilibrium. In contemporary communication architectures, however, dark-social cascades accelerate during the final 48 hours of campaigning, precisely when statutory blackout periods and methodological field cutoffs suspend polling updates. Bridging Converse's belief constraint paradigm with Zaller's Receive-Accept-Sample model, the two-timescale formulation couples opinions to convictions through an elastic hyperbolic-tangent tether. Transient campaign blitzes displace expressed opinions $\mathbf{x}_i$ away from internal baselines, while slow-crystallizing convictions $\mathbf{s}_i$ exert a restorative drag that pulls opinions back once informational shocks subside. Consequently, an acute campaign blitz can produce sharp short-term voting swings that alter ballot outcomes without requiring deep ideological conversion, explaining how pre-election surveys can miss late-breaking inflection points.

My results suggest that this late-campaign electoral vulnerability is fundamentally governed by cognitive negativity asymmetry. In the calibrated model, candidate aversion accumulates with rapid sensitivity while decaying with extreme hysteresis ($\lambda_R = 0.14, \delta_R = 0.04$), producing a kinetic acquisition-to-decay ratio of $\lambda_R / \delta_R = 3.5\times$. In contrast, positive candidate admiration accumulates gradually and decays rapidly ($\lambda_A = 0.12, \delta_A = 0.06$). When combined with voting decision utilities that penalize candidate aversion more heavily than they reward positive affinity ($\theta_{\mathrm{rej}} = 1.65$), this kinetic asymmetry confers an inherent structural advantage upon negative disinformation blitzes. Positive candidate enthusiasm requires weeks of sustained communicative cultivation; in contrast, visceral candidate aversion can be synthesized within a 48-hour window, creating psychological scars that linger well beyond election day.

Furthermore, the multi-layer architecture exposes a structural pathology in contemporary platform curation. Commercial recommendation algorithms optimize for gross user engagement by weighting affirmative reactions and hostile outrage symmetrically, rewarding moral-emotional controversy along the diagonal where partisan admiration meets out-group outrage. Under this curation logic, conventional democratic defenses systematically misfire. Institutional newsrooms and regulatory courts operate under a wide separation of operational timescales ($\tau_{\mathrm{turnaround}} \gg \tau_{\mathrm{social}}$). In practice, defensive counter-speech, outrage debunking, and adversarial quote-tweets supply the out-group friction signals that recommendation engines use to distribute opponent messages into unaligned discovery feeds.

The second-round turnaround suggests that democratic resilience in algorithmic environments benefits from active socio-technical engagement rather than passive legalism alone. By translating the operational doctrine of \emph{Janonismo Cultural} into general graph-theoretic mechanisms, the model illustrates how coordinated socio-technical counter-strategies can systematically neutralize engagement-driven offensives. An outrage-starvation embargo (T1) deprives recommendation engines of engagement telemetry, collapsing opponent reach by $48\%$. This observation resonates with recent agent-based analyses of strategic behavioral norms in peer-to-peer networks, which indicate that simple shifts in local interaction rules and cooperative or defensive norms, such as withholding reactive engagement rather than amplifying friction, can nonlinearly transform macroscopic collective outcomes \citep{alvarenga2026power}. In parallel, monothematic bandwidth concentration (T2) exploits superlinear returns in platform discovery functions, quadrupling viral reach over dispersed messaging. Loss-aversion framing (T3) activates Kahneman--Tversky loss weighting (a 2.25-fold aversion to losses relative to equivalent gains), mobilizing cross-pressured working-class swing voters. Finally, molecular encrypted cascades (T4), physical turnout locking (T5), and attention-ceiling saturation (T6) bypass open-feed algorithms, insulate voters from late-stage blitzes, and cross the majority threshold.

\subsection{Platform Governance and Cross-Border Spillovers}
\label{subsec:platform_governance}
Beyond domestic electoral competition, these platform dynamics carry critical geopolitical and regulatory implications for socio-technical governance and transnational information flows. When the primary deliberative infrastructure of modern democracies relies on proprietary recommendation engines controlled by offshore technology corporations outside national jurisdiction, electoral integrity becomes structurally vulnerable to opaque curation weights and corporate content moderation policies. During recent Brazilian electoral cycles, major digital platforms faced public scrutiny, parliamentary inquiries, and regulatory compliance orders from the Superior Electoral Court regarding dark-money advertising reach and algorithmic virality disparities. 

Furthermore, because the scale-free topology of global social media platforms aggregates user engagement signals across jurisdictional borders, synchronized domestic networks can trigger transnational algorithmic spillovers. This phenomenon was empirically demonstrated when the 10-million-member \emph{Ex\'ercito do Lula} messaging collective redirected its coordinated mobilization infrastructure toward amplifying United States political campaign narratives. Because platform recommendation algorithms weigh high-velocity engagement cascades without intrinsic regard for geographic origin, large-scale coordinated domestic interactions can shift global relevance metrics ($R_m$) and breach algorithmic distribution thresholds across borders. These cross-border interactions highlight the challenge of reconciling platform auditing with national democratic oversight.

\subsection{Modeling Boundaries and Future Extensions}
\label{subsec:limitations}
Several modeling boundaries should be highlighted. First, the agent-based simulations treat interpersonal multiplex ties and platform follower networks as stationary topologies over the 50-day campaign horizon. While empirical network rewiring is negligible over weeks, multi-year electoral polarization inevitably induces selective tie severance and echo-chamber hardening. Second, the framework models political attitudes within a continuous metric concept hypercube $\Omega = [-1, 1]^4$. Although this captures core policy dimensions with high quantitative fidelity, it abstracts away linguistic nuances, narrative irony, and generative AI synthetic media. Future work should integrate large language model (LLM) agents capable of generating and interpreting unconstrained multimodal political rhetoric. Third, the second-round counterfactuals assume bounded counter-adaptation by political opponents to isolate short-term tactical shocks. In iterative electoral cycles, rival campaigns will adapt their own playbooks, sparking dynamic co-evolutionary arms races.

\section{Conclusion}
\label{sec:conclusion}

Despite these modeling simplifications, the core behavioral mechanisms identified in this study, the two-timescale opinion-conviction tether, cognitive negativity hysteresis, and symmetric outrage curation, represent generalizable features of contemporary information ecosystems. Addressing democratic vulnerability requires structural platform reforms, such as deprecating out-group hostility signals in recommendation rankings, paired with the strategic capacity of democratic campaigns to defend popular majorities within algorithmic environments. By demonstrating how the same fast-opinion elasticity that enables devastating election-eve disinformation blitzes can also be leveraged through disciplined, socio-technical counter-strategies, this study establishes a computational foundation for understanding and engineering electoral resilience in the digital age.

\section*{CRediT authorship contribution statement}
\textbf{César Rennó-Costa}: Conceptualization, Methodology, Software, Validation, Formal analysis, Investigation, Data curation, Writing - original draft, Writing - review \& editing, Visualization, Supervision, Project administration, Funding acquisition.

\section*{Declaration of competing interest}
The author declares no known competing financial interests or personal relationships that could have influenced the work reported in this paper.

\section*{Data and code availability}
The empirical demographic, polling, broadcast airtime, and election result datasets supporting the findings of this study, alongside all simulation code, Numba JIT kernels, graph generation scripts, and figure plotting pipelines, are archived and openly available under the MIT License at \url{https://github.com/cesarrennocosta/janones}.

\section*{Declaration of generative AI and AI-assisted technologies in the writing process}
During the preparation of this work the author used Claude and Gemini in order to assist with code optimization, data formatting, and language polishing. After using these tools, the author reviewed and edited the content as needed and takes full responsibility for the content of the published article.

\section*{Acknowledgements}
I thank Luana Alvarenga, Renata Mota, and Fernando Pureza for insightful discussions and feedback. I acknowledge the high-performance computing resources provided by the N{\'u}cleo de Processamento de Alto Desempenho (NPAD) at the Federal University of Rio Grande do Norte (UFRN). This work received financial and institutional support from the Coordena{\c{c}}{\~a}o de Aperfei{\c{c}}oamento de Pessoal de N{\'i}vel Superior (CAPES), the Conselho Nacional de Desenvolvimento Cient{\'i}fico e Tecnol{\'o}gico (CNPq), and the Pr{\'o}-Reitoria de Pesquisa of the Federal University of Rio Grande do Norte (UFRN/PROPESQ).

\appendix
\section{Model Parameters and Tactical Specifications}
\label{app:parameters}

Tables~\ref{tab:supp_model_parameters}, \ref{tab:supp_media_parameters}, and \ref{tab:supp_janones_tactics} provide the comprehensive parameterization and tactical mapping of the multi-layer digital twin.

\begin{table*}[t!]
\centering
\caption{Summary of population demographics, interpersonal multiplex network, and two-timescale conviction--opinion parameters.}
\label{tab:supp_model_parameters}
\footnotesize
\begin{tabular*}{\textwidth}{@{\extracolsep{\fill}}lp{5.2cm}p{3.8cm}p{4.2cm}@{}}
\toprule
\textbf{Symbol} & \textbf{Description} & \textbf{Value / Range} & \textbf{Reference} \\
\midrule
\multicolumn{4}{l}{\textit{Global Simulation \& Concept Space}} \\
$N$ & City population size & $\{10^4, 10^5, 10^6, 10^7\}$ & \cite{bettencourt2007growth} \\
$D$ & Number of policy/concept dimensions & $4$ ($\Omega = [-1, 1]^4$) & \cite{converse1964nature} \\
$T, \Delta t$ & Horizon and integration step size & $T \in \{100, 1500\}$, $\Delta t = 0.15$ & \cite{friedkin1990social} \\
$S$ & Monte Carlo realizations per scenario & $100$ & This study \\
$\mathbf{C}$ & Cognitive issue-coupling matrix & $\mathbb{R}^{4 \times 4}$, Eq.~\eqref{eq:coupling_matrix} & \cite{friedkin2016network} \\
\midrule
\multicolumn{4}{l}{\textit{Citizen Demographics \& Psychological Traits}} \\
$a_i, \tilde{a}_i$ & Citizen age (years) and normalized age & $[16, 90]$, $\tilde{a}_i \in [0, 1]$ & \cite{ibge2022censo} \\
$\mathrm{SES}_i$ & Socioeconomic and educational index & $\mathrm{Beta}(2, 3) \in [0, 1]$ & \cite{ibge2022censo} \\
$\mathbf{s}_i(t)$ & Deep conviction vector of citizen $i$ & $[-1, 1]^D$ & \cite{converse1964nature, kiley2020measuring} \\
$\mathbf{x}_i(t)$ & Current opinion vector of citizen $i$ & $[-1, 1]^D$ & \cite{zaller1992nature} \\
$E_i(t)$ & Conviction extremity ($\|\mathbf{s}_i(t)\|_2 / \sqrt{D}$) & $[0, 1]$ & \cite{howe2017attitude} \\
$\alpha_i$ & Interpersonal social susceptibility & $[0.05, 0.98]$ & \cite{friedkin1990social, alwin1991aging} \\
$\epsilon_i$ & Bounded confidence (base $\bar{\epsilon}=0.45$) & $[0.10, 1.50]$ & \cite{deffuant2000mixing, hegselmann2002opinion} \\
$u_i(t)$ & Attitude uncertainty (Relative Agreement) & $[0.02, 1.00]$ & \cite{deffuant2002how} \\
$b_i$ & Biased assimilation (base $\bar{b}=1.25$) & $[0.00, 5.00]$ & \cite{lord1979biased, dandekar2013biased} \\
\midrule
\multicolumn{4}{l}{\textit{Five-Layer Multiplex Interpersonal Network}} \\
$w_1, \langle s_h \rangle$ & Layer 1 (\texttt{family}): household clique & $w_1 = 0.36$, $\langle s_h \rangle = 2.79$ & \cite{ibge2022censo} \\
$w_2, \langle k^{[2]} \rangle$ & Layer 2 (\texttt{work\_school}): peer degree & $w_2 = 0.26$, $\langle k^{[2]} \rangle = 10$ & \cite{inep2023censo, mte2023rais} \\
$w_3, \langle k^{[3]} \rangle$ & Layer 3 (\texttt{church\_community}): degree & $w_3 = 0.14$, $\langle k^{[3]} \rangle = 8$ & \cite{mcpherson2001birds} \\
$w_4, \langle k^{[4]} \rangle, p_{\mathrm{ws}}$ & Layer 4 (\texttt{friends}): small-world circle & $w_4 = 0.14$, $\langle k^{[4]} \rangle = 6$, $p_{\mathrm{ws}}=0.15$ & \cite{watts1998collective} \\
$w_5, \langle k^{[5]} \rangle$ & Layer 5 (\texttt{weak\_ties}): transit ties & $w_5 = 0.10$, $\langle k^{[5]} \rangle = 4$ & \cite{granovetter1973strength, ibge2022censo} \\
\midrule
\multicolumn{4}{l}{\textit{Two-Timescale Conviction--Opinion \& Fact Parameters}} \\
$\Delta_0, \Delta_{\max, i}$ & Base and individual opinion tether radius & $\Delta_0 = 0.34$, Eq.~\eqref{eq:methods_tether_radius} & \cite{sherif1961social, howe2017attitude} \\
$\eta_0$ & Base conviction adaptation drift rate & $0.0024 \text{ step}^{-1}$ & \cite{kiley2020measuring} \\
$\tau_c$ & Conviction crystallization timescale & $280 \text{ steps}$ & \cite{alwin1991aging} \\
$\gamma_0, \sigma_{\xi}$ & Factual pull and cognitive perturbation noise & $\gamma_0 = 0.055 \text{ step}^{-1}$, $\sigma_{\xi} = 0.012$ & This study \\
\bottomrule
\end{tabular*}
\end{table*}

\begin{table*}[t!]
\centering
\caption{Summary of traditional mass media and algorithmic social media platform parameters.}
\label{tab:supp_media_parameters}
\footnotesize
\begin{tabular*}{\textwidth}{@{\extracolsep{\fill}}lp{5.2cm}p{3.8cm}p{4.2cm}@{}}
\toprule
\textbf{Symbol} & \textbf{Description} & \textbf{Value / Range} & \textbf{Reference} \\
\midrule
\multicolumn{4}{l}{\textit{Traditional Mass Media Ecosystem}} \\
$M, \mathbf{m}_g, \tau_g$ & Media groups, ideological centers, tendentiousness & $M \in \{0, 1, 2, 3, 4, 12\}, \tau_g \in [0, 1]$ & \cite{mullainathan2005market, gentzkow2006media} \\
$v \in \{0, 1, 2, 3, 4\}$ & Modalities (TV, Radio, Press, Outdoor, Rally) & 5 channels per conglomerate & \cite{kantar2023inside, reuters2023digital} \\
$\gamma_v, \delta_v$ & Modality credibility and attention decay rate & $\gamma \in [0.45, 0.82], \delta \in [0.08, 0.70]$ & \cite{wu2007novelty, downs1972up} \\
$\beta_m, \kappa_0$ & Media selective exposure beta and broadcast strength & $\beta_m \in [0.15, 2.16], \kappa_0 = 0.05$ & \cite{dandekar2013biased, dellavigna2007fox} \\
\midrule
\multicolumn{4}{l}{\textit{Algorithmic Social Media Platforms \& Influencers}} \\
$p \in \{0, 1, 2, 3\}$ & Platforms (Instagram/TikTok, YouTube/X, FB, WhatsApp) & 4 platform modalities & \cite{reuters2023digital, cinelli2021echo} \\
$\gamma_{\mathrm{pl}}$ & Influencer follower power-law exponent (Gini $\approx 0.73$) & $\gamma_{\mathrm{pl}} \approx 1.15$ & \cite{barabasi1999emergence, bianconi2001competition} \\
$\beta_{\mathrm{cur}}, \mu_{\mathrm{emo}}$ & Algorithmic homophily and creator emotionality weight & $\beta_{\mathrm{cur}} = 2.0, \mu_{\mathrm{emo}} = 1.5$ & \cite{brady2017emotion, vosoughi2018spread} \\
$\omega$ & Dual-valence engagement weight ($R_m$) & $\omega = 1.0$ & \cite{crockett2017moral, bail2018exposure} \\
$\tau_{\text{WA}}$ & WhatsApp peer trust weight ($K \le 256$) & $\tau_{\text{WA}} = 0.85$ & \cite{resende2019whatsapp, garimella2020images} \\
$\tau_{\mathrm{inf}}$ & Influencer creator tendentiousness & $[0.10, 0.90]$ & \cite{sasahara2021social, tucker2018social} \\
\midrule
\multicolumn{4}{l}{\textit{Electoral Decision Calculus \& Asymmetric Kinetics}} \\
$\lambda_A, \delta_A$ & Candidate admiration acquisition and decay & $\lambda_A = 0.12, \delta_A = 0.06$ & This study \\
$\lambda_R, \delta_R$ & Candidate aversion acquisition and decay & $\lambda_R = 0.14, \delta_R = 0.04$ ($\lambda_R/\delta_R = 3.5\times$) & \cite{baumeister2001bad, rozin2001negativity} \\
$\theta_{\mathrm{rej}}$ & Quadratic aversion penalty in decision utility & $\theta_{\mathrm{rej}} = 1.65$ & \cite{almeida2019evangelicals} \\
$\beta_{\mathrm{vote}}, \tau_{\mathrm{abst}}$ & Logit choice sensitivity and base abstention threshold & $\beta_{\mathrm{vote}} = 3.2, \tau_{\mathrm{abst}} = 0.38$ & \cite{mcfadden1974conditional} \\
$\mu_{\mathrm{turn}}, \delta_{\mathrm{transit}}$ & Turnout logistic steepness and transit subsidy & $\mu_{\mathrm{turn}} = 4.5, \delta_{\mathrm{transit}} = -0.15$ & This study \\
\bottomrule
\end{tabular*}
\end{table*}

\begin{table*}[t!]
\centering
\caption{Tactical translation table: operationalizing the Janones doctrine into agent-based graph mechanics across the foundational model layers.}
\label{tab:supp_janones_tactics}
\footnotesize
\begin{tabular*}{\textwidth}{@{\extracolsep{\fill}}lp{4.2cm}p{5.0cm}p{3.6cm}@{}}
\toprule
\textbf{Janones Tactic} & \textbf{Operational Description} & \textbf{Agent / Graph Mechanism} & \textbf{Parameter / Subsystem} \\
\midrule
T1 Asphyxiation & Suppression of outbound interaction toward rival bait & Interactivity probability $P_{\mathrm{reply}}, P_{\mathrm{quote}} \to 0$ & $R_m$ signal denial \\
T2 Monothematism & One coordinated agenda per 24h cycle & Variance $\sigma^2(T_{\mathrm{day}}) \to 0$ across allied nodes & Quadratic virality $\sum x_k^2$ \\
T3 Material Loss Frame & Concrete economic risks (welfare cuts, wage freeze) & Kahneman--Tversky loss weighting & $\lambda_{\mathrm{loss}} \in [2.0, 2.5]$ \\
T4 WhatsApp Sync & Scripted asset dispatch and molecular relay & Phased dispatch into private cliques & Trust weight $\times 2.8$ ($\tau_{\text{WA}} = 0.85$) \\
T5 Mass Rallies & Vast public assemblies in metropolitan capitals & Modality $v=4$ and multiplex layers $\mathcal{L}_1/\mathcal{L}_4$ & Turnout locking $A_{i,L} > \tau_{\mathrm{turnout}}$ \\
T6 24h Live Saturation & Nonstop continuous broadcast during pre-election sprint & Attention saturation ceiling & Capping rival blitz by $65\%$ \\
\bottomrule
\end{tabular*}
\end{table*}

\bibliographystyle{cas-model2-names}
\bibliography{references}

\end{document}
